\chapter[Homogenizacja okresowa równań eliptycznych]{Homogenizacja równań eliptycznych: zadanie komórkowe, macierz efektywna i granica oscylacyjna}
\label{chap:pde30}

\section{Dlaczego współczynników nie można po prostu uśrednić?}
Rozważmy ograniczony obszar Lipschitzowski $\Omega\subset\R^d$, $d\ge1$,
liczbę $\eps>0$ i macierz $A(y)$, która zmienia się w skali
pojedynczej komórki $Y=(0,1)^d$. W zagadnieniu Dirichleta
\begin{equation}\label{eq:pde30-micro}
 -\diver\bigl(A(x/\eps)\nabla u_\eps\bigr)=f\quad\text{w }\Omega,
 \qquad u_\eps\in H_0^1(\Omega),
\end{equation}
rozwiązanie odczuwa coraz drobniejsze niejednorodności. Pytamy,
czy istnieje granica $u_0$ i czy spełnia równanie tej samej postaci,
z pewną stałą macierzą $A^{\mathrm{hom}}$. Samo zastąpienie $A$
jego średnią arytmetyczną jest na ogół błędne: w kierunku
prostopadłym do warstw o przepływie decyduje średnia harmoniczna.

Przyjmujemy okres komórki równy $1$ w każdym kierunku.
Jest to inna normalizacja geometryczna niż dla torusa
$\T^d=\R^d/(2\pi\mathbb Z)^d$ z rozdziału~\ref{chap:pde25};
przy porównywaniu wzorów trzeba przeskalować zmienne.
Pracujemy na rzeczywistych przestrzeniach funkcyjnych.
Przez $H^{-1}(\Omega)$ rozumiemy przestrzeń dualną
do $H_0^1(\Omega)$ z normą względem normy gradientu
$\|v\|_{H_0^1}=\|\nabla v\|_2$,
równoważnej normie Sobolewa na mocy
Wniosku~\ref{cor:pde10-Poincare}.
Przez $H^1_{\mathrm{per}}(Y)$ rozumiemy przestrzeń ograniczeń do
$Y$ funkcji z $H^1_{\mathrm{loc}}(\R^d)$, których przedłużenie jest
$\mathbb Z^d$-okresowe. Jej podprzestrzenią o średniej zero jest
$H^1_{\mathrm{per},0}(Y)=\{v\in H^1_{\mathrm{per}}(Y):
\int_Yv\,\mathrm d\lambda_d(y)=0\}$.
Niech $A\in L^\infty_{\mathrm{per}}(Y;\R^{d\times d})$ będzie
\emph{symetryczna} oraz niech $0<\lambda\le\Lambda$ spełniają
\begin{equation}\label{eq:pde30-ellipticity}
 \lambda|\xi|^2\le\xi^{\mathsf T}A(y)\xi\le\Lambda|\xi|^2
 \quad (\xi\in\R^d),\quad\text{dla prawie każdego }y\in Y.
\end{equation}
Macierz przedłużamy okresowo na $\R^d$.
Przeskalowanie $x\mapsto x/\eps$ zachowuje zbiory
miary zero w tym sensie, że ich przeciwobrazy
również mają miarę zero. Z symetrii i warunku
~\eqref{eq:pde30-ellipticity} wynika
$\|A(y)\|_{\mathrm{op}}\le\Lambda$ prawie wszędzie.
Definiujemy na $H_0^1(\Omega)$ formę
\[
 a_\eps(u,v)=\int_\Omega A(x/\eps)\nabla u\cdot\nabla v\,
 \mathrm d\lambda_d(x).
\]
Mamy $|a_\eps(u,v)|\le\Lambda\|\nabla u\|_2\|\nabla v\|_2$
oraz $a_\eps(v,v)\ge\lambda\|\nabla v\|_2^2$.
Z twierdzenia Laxa--Milgrama dla każdego $f\in H^{-1}(\Omega)$
istnieje dokładnie jedno $u_\eps\in H_0^1(\Omega)$
spełniające $a_\eps(u_\eps,v)=\langle f,v\rangle$
dla każdego $v\in H_0^1(\Omega)$.
Jest to równoważne równaniu~\eqref{eq:pde30-micro}
w sensie dystrybucji przez gęstość funkcji testowych.
Testując przez $v=u_\eps$, dostajemy
$\lambda\|\nabla u_\eps\|_2^2\le
\|f\|_{H^{-1}}\|\nabla u_\eps\|_2$,
a stąd $\|\nabla u_\eps\|_2\le\lambda^{-1}\|f\|_{H^{-1}}$
(również gdy gradient znika).
Zapis $\diver$ oznacza dywergencję, a $\mathrm d\lambda_d$ miarę Lebesgue'a.

\section{Przypadek jednowymiarowy: dokładny wzór i tempo zbieżności}
Niech $\Omega=(0,L)$, $L>0$, $a\in L^\infty_{\mathrm{per}}(0,1)$
oraz $0<\lambda\le a\le\Lambda$ prawie wszędzie. Definiujemy
$b=1/a$ i $\bar b=\int_0^1b(y)\,\mathrm d\lambda_1(y)$, a także
$a_{\mathrm{hom}}=\bar b^{-1}$. Dla $f\in L^2(0,L)$ połóżmy
$F(x)=\int_0^xf(t)\,\mathrm d\lambda_1(t)$.
Mamy $1/\Lambda\le b\le1/\lambda$, więc
$\lambda\le a_{\mathrm{hom}}\le\Lambda$.
Funkcja $F$ należy do $H^1(0,L)$ i ma ciągły,
absolutnie ciągły reprezentant, $F'=f$ prawie wszędzie,
oraz $\|F\|_\infty\le\sqrt L\|f\|_2$.

\begin{twierdzenie}[Jednowymiarowe równanie uśrednione]\label{thm:pde30-1d}
Dla każdego $\eps>0$ istnieje dokładnie jedno $u_\eps\in H_0^1(0,L)$
spełniające $-(a(x/\eps)u_\eps')'=f$ w sensie dystrybucji.
Niech $u_0$ będzie jedynym rozwiązaniem
$-a_{\mathrm{hom}}u_0''=f$ w $H_0^1(0,L)$.
Wtedy, gdy $\eps\downarrow0$,
$u_\eps\to u_0$ w $L^\infty(0,L)$ oraz
$u_\eps\rightharpoonup u_0$ w $H_0^1(0,L)$, a strumienie
$a(x/\eps)u_\eps'$ zbiegają silnie w $L^2(0,L)$ do
$a_{\mathrm{hom}}u_0'$. Istnieje stała $C$, zależna tylko od
$L,\lambda,\Lambda$, taka że dla $0<\eps\le1$
\begin{equation}\label{eq:pde30-1d-rate}
 \|u_\eps-u_0\|_{L^\infty(0,L)}
 +\|a(x/\eps)u_\eps'-a_{\mathrm{hom}}u_0'\|_{L^2(0,L)}
 \le C\eps\|f\|_{L^2(0,L)}.
\end{equation}
\end{twierdzenie}
\begin{proof}
Istnienie i jednoznaczność wynikają z poprzedniej sekcji.
Ponieważ $q_\eps=a(x/\eps)u_\eps'\in L^2(0,L)$
i $q_\eps'=-f\in L^2(0,L)$ w sensie dystrybucji,
mamy $q_\eps\in H^1(0,L)$. Z jednowymiarowej
charakterystyki Sobolewa (Twierdzenie~\ref{thm:pde9-ac})
otrzymujemy absolutnie ciągły reprezentant
$q_\eps(x)=C_\eps-F(x)$. Całkując $u_\eps'$ oraz
wykorzystując $u_\eps(L)=u_\eps(0)=0$, otrzymujemy
\begin{equation}\label{eq:pde30-explicit-1d}
 C_\eps=\frac{\int_0^L b(t/\eps)F(t)\,\mathrm d\lambda_1(t)}
 {\int_0^L b(t/\eps)\,\mathrm d\lambda_1(t)},\qquad
 u_\eps(x)=\int_0^xb(t/\eps)(C_\eps-F(t))\,\mathrm d\lambda_1(t).
\end{equation}
Odwrotnie, wzór~\eqref{eq:pde30-explicit-1d}
określa funkcję absolutnie ciągłą z pochodną
$b(x/\eps)(C_\eps-F(x))\in L^2$, zerowymi wartościami
na końcach i strumieniem $C_\eps-F$; jest zatem
rozwiązaniem zagadnienia. W rozwiązaniu granicznym
analogicznie $u_0'=\bar b(C_0-F)$,
gdzie $C_0=L^{-1}\int_0^LF(t)\,\mathrm d\lambda_1(t)$.
W szczególności $u_0\in H^2(0,L)$ i
$u_0''=-\bar b f$.

Wyjaśnijmy ilościowo uśrednianie. Jeśli $g\in L^\infty_{\mathrm{per}}(0,1)$
ma średnią zero, to jej funkcja pierwotna
$G(s)=\int_0^sg(y)\,\mathrm d\lambda_1(y)$ jest ograniczona
i okresowa, ponieważ
$G(s+1)-G(s)=\int_s^{s+1}g\,\mathrm d\lambda_1=0$.
Ponadto $\|G\|_\infty\le\|g\|_{L^1(0,1)}$.
Dla $h\in W^{1,1}(0,L)$ stosujemy regułę iloczynu
dla absolutnie ciągłych reprezentantów oraz
$(\eps G(x/\eps))'=g(x/\eps)$ prawie wszędzie:
\[
 \int_0^Lg(x/\eps)h(x)\,\mathrm d\lambda_1(x)
 =\eps[G(x/\eps)h(x)]_0^L
 -\eps\int_0^LG(x/\eps)h'(x)\,\mathrm d\lambda_1(x).
\]
Stąd
\begin{equation}\label{eq:pde30-average-rate}
 \left|\int_0^Lg(x/\eps)h(x)\,\mathrm d\lambda_1(x)\right|
 \le\eps\|G\|_\infty
 \bigl(|h(0)|+|h(L)|+\|h'\|_{L^1(0,L)}\bigr).
\end{equation}
Połóżmy $g=b-\bar b$ oraz $h=F-C_0$.
Mamy $\int_0^L h\,\mathrm d\lambda_1=0$,
$|C_0|\le\|F\|_\infty$ i
$\|h\|_\infty\le2\sqrt L\|f\|_2$,
$\|h'\|_1\le\sqrt L\|f\|_2$.
Odejmując $C_0$ od ilorazu
~\eqref{eq:pde30-explicit-1d}, dostajemy
\[
 C_\eps-C_0=
 \frac{\int_0^L g(t/\eps)h(t)\,\mathrm d\lambda_1(t)}
 {\int_0^Lb(t/\eps)\,\mathrm d\lambda_1(t)}.
\]
Mianownik jest nie mniejszy od $L/\Lambda$,
więc~\eqref{eq:pde30-average-rate} daje
$|C_\eps-C_0|\le C\eps\|f\|_2$.
Stała zależy wyłącznie od $L,\lambda,\Lambda$,
gdyż $\|G\|_\infty\le\|b-\bar b\|_1\le2/\lambda$.
Następnie odejmujemy wzory na pochodne i całkujemy:
\[
 \begin{aligned}
 u_\eps(x)-u_0(x)
 &=(C_\eps-C_0)\int_0^x b(t/\eps)\,\mathrm d\lambda_1(t)\\
 &\quad+\int_0^x g(t/\eps)(C_0-F(t))\,\mathrm d\lambda_1(t).
 \end{aligned}
\]
Pierwszy składnik jest ograniczony przez
$(L/\lambda)|C_\eps-C_0|$, a drugi przez
$\eps\|G\|_\infty(2\|h\|_\infty+\|h'\|_1)$
na mocy~\eqref{eq:pde30-average-rate} na $(0,x)$.
Otrzymujemy oszacowanie jednostajne na $[0,L]$.
Ponieważ $q_\eps-q_0=C_\eps-C_0$, norma
$L^2$ różnicy strumieni jest równa
$\sqrt L|C_\eps-C_0|$; w istocie zachodzi
także zbieżność strumieni w normie jednostajnej.
To dowodzi~\eqref{eq:pde30-1d-rate}.
Jednostajne ograniczenie gradientów i refleksywność
$H_0^1$ pozwalają z dowolnego ciągu $\eps\to0$
wydzielić słabo zbieżny podciąg. Ciągłe zanurzenie
w $L^2$ i wykazana zbieżność jednostajna identyfikują
jego słabą granicę z $u_0$. Ponieważ każdy ciąg
ma taki podciąg, cała rodzina zbiega słabo do $u_0$
(testując przez dowolny ciągły funkcjonał liniowy).
\end{proof}

\begin{stwierdzenie}[Poprawka do gradientu w wymiarze jeden]\label{prop:pde30-corrector1d}
Przy założeniach twierdzenia~\ref{thm:pde30-1d} istnieje jedyna
funkcja $\chi\in W^{1,\infty}_{\mathrm{per}}(0,1)$ o średniej zero,
której pochodna spełnia $\chi'(y)=a_{\mathrm{hom}}/a(y)-1$.
Ponadto
\[
\left\|u_\eps'-\bigl(1+\chi'(x/\eps)\bigr)u_0'\right\|_{L^2(0,L)}
\le C\eps\|f\|_2.
\]
Jeśli $a$ nie jest prawie wszędzie stała i $f\not\equiv0$,
to nie zachodzi silna zbieżność $u_\eps'\to u_0'$.
\end{stwierdzenie}
\begin{proof}
Funkcja $g_1=a_{\mathrm{hom}}b-1$ należy do
$L^\infty_{\mathrm{per}}(0,1)$ i ma średnią zero,
bo $a_{\mathrm{hom}}\bar b=1$.
Jej funkcja pierwotna $G_1(y)=\int_0^y g_1\,\mathrm d\lambda_1$
jest okresowa i lipschitzowska. Definiujemy
$\chi=G_1-\int_0^1G_1\,\mathrm d\lambda_1$.
Dwie funkcje o wymaganej pochodnej różnią się stałą,
a średnia zero wymusza jej znikanie.
Z tożsamości $u_0'=\bar b(C_0-F)$ wynika
$(1+\chi'(x/\eps))u_0'=b(x/\eps)(C_0-F)$.
Różnica z $u_\eps'$ to $b(x/\eps)(C_\eps-C_0)$,
więc twierdzenie~\ref{thm:pde30-1d} daje oszacowanie.
Jeśli $a$ nie jest stała, to średnia z
$(a_{\mathrm{hom}}/a-1)^2$ jest dodatnia.
Przy $f\not\equiv0$ mamy $u_0'\not\equiv0$,
bo $u_0''=-\bar b f$.
Ponieważ $u_0'\in H^1(0,L)\subset L^\infty(0,L)$,
waga $|u_0'|^2$ należy do $W^{1,1}(0,L)$
z pochodną $2u_0'u_0''$.
Stosujemy~\eqref{eq:pde30-average-rate} do
$g_1^2-\int_0^1g_1^2\,\mathrm d\lambda_1$
oraz tej wagi. Różnica
$u_\eps'-u_0'-g_1(x/\eps)u_0'$ dąży do zera w $L^2$
z już wykazanego oszacowania poprawki,
a normy $g_1(x/\eps)u_0'$ są jednostajnie ograniczone.
Zatem różnica kwadratów norm również dąży do zera,
co daje
\[
\lim_{\eps\downarrow0}\|u_\eps'-u_0'\|_2^2=
\left(\int_0^1|a_{\mathrm{hom}}/a(y)-1|^2
\,\mathrm d\lambda_1(y)\right)\|u_0'\|_2^2>0.
\]
Istotne są więc poprawki oscylacyjne, a nie tylko granica $u_0$.
\end{proof}

\begin{przyklad}[Dlaczego średnia arytmetyczna jest błędna]
Niech na połowie komórki $a=1$, a na drugiej połowie $a=4$.
Wówczas średnia arytmetyczna wynosi $5/2$, a efektywny
współczynnik to $a_{\mathrm{hom}}=(\tfrac12+\tfrac18)^{-1}=8/5$.
Dla $f=1$ na $(0,1)$ rozwiązanie graniczne ma postać
$u_0(x)=x(1-x)/(2a_{\mathrm{hom}})$.
Zastąpienie $a_{\mathrm{hom}}$ przez $5/2$ daje
$u_{\mathrm{ar}}(x)=x(1-x)/5$ zamiast
$u_0(x)=5x(1-x)/16$. Różnica jest równa
$9x(1-x)/80$ i ma normę jednostajną $9/320$;
nie znika mimo dowolnie drobnych warstw.
Dla tych dwóch faz $g_1$ przyjmuje wartości
$3/5$ i $-3/5$, więc
$\int_0^1g_1^2\,\mathrm d\lambda_1=9/25$.
Ponieważ $u_0'(x)=5(1-2x)/16$,
$\|u_0'\|_2^2=25/768$. Ze Stwierdzenia
~\ref{prop:pde30-corrector1d} otrzymujemy
$\lim_{\eps\downarrow0}\|u_\eps'-u_0'\|_2^2=3/256$.
\end{przyklad}

\section{Uśrednianie funkcji okresowych}
Do przejścia w wielu wymiarach potrzebujemy prostego lematu.
\begin{lemat}[Słaba granica funkcji okresowych]\label{lem:pde30-average}
Niech $D\subset\R^d$ będzie ograniczonym zbiorem otwartym,
a $g\in L^2_{\mathrm{per}}(Y)$ (czyli $g$ jest okresowym
przedłużeniem funkcji z $L^2(Y)$). Wtedy
$g(x/\eps)\rightharpoonup\bar g:=\int_Yg(y)\,\mathrm d\lambda_d(y)$
słabo w $L^2(D)$, gdy $\eps\downarrow0$.
Ponadto $\eps g(x/\eps)\to0$ silnie w $L^2(D)$.
\end{lemat}
\begin{proof}
Wybierzmy $R\ge1$ takie, że $D\subset(-R,R)^d$.
Dla $0<\eps\le1$ co najwyżej $(2R/\eps+2)^d$
komórek $\eps(k+Y)$, $k\in\mathbb Z^d$, przecina
ten sześcian. Ich brzegi mają miarę zero.
Przez podstawienie $x=\eps(k+y)$ i okresowość
całka na każdej pełnej komórce wynosi
$\eps^d\|g\|_{L^2(Y)}^2$. W szczególności
\[
 \|g(\cdot/\eps)\|_{L^2(D)}
 \le (2R+2)^{d/2}\|g\|_{L^2(Y)}.
\]
To daje również silną zbieżność $\eps g(\cdot/\eps)$
do zera w $L^2(D)$.

Wykażemy słabą zbieżność bez założeń o regularności brzegu.
Niech $h=g-\bar g$ i $\varphi\in C_c^\infty(D)$,
przedłużona przez zero na $\R^d$. Z zerowej średniej
$h$ na każdej przeskalowanej komórce wynika
\[
 \begin{aligned}
 \int_D h(x/\eps)\varphi(x)\,\mathrm d\lambda_d(x)
 =\sum_{k\in\mathbb Z^d}\eps^d\int_Y h(y)
 \bigl(\varphi(\eps(k+y))-\varphi(\eps k)\bigr)
 \,\mathrm d\lambda_d(y).
 \end{aligned}
\]
Wystarczy sumować po komórkach, których domknięcia
przecinają nośnik $\varphi$; pozostałe składniki są zerowe.
Ich liczba jest nie większa niż $(2R/\eps+2)^d$.
Różnica wartości funkcji testowej jest ograniczona
przez $\eps\sqrt d\|\nabla\varphi\|_\infty$.
Moduł całej sumy jest więc nie większy niż
\[
 (2R+2)^d\eps\sqrt d\|\nabla\varphi\|_\infty
 \|h\|_{L^1(Y)},
\]
co dąży do zera. Dla dowolnego testu $\varphi\in L^2(D)$
wybieramy $\varphi_n\in C_c^\infty(D)$ zbieżne
do niego w $L^2$. Błąd zastąpienia testu przez
$\varphi_n$ jest ograniczony przez
$(2R+2)^{d/2}\|h\|_{L^2(Y)}\|\varphi-\varphi_n\|_2$,
jednostajnie względem $\eps$. Najpierw przechodzimy
do granicy z $\eps$, potem z $n$.
Otrzymujemy tezę o słabej zbieżności.
\end{proof}

\section{Zadanie komórkowe i macierz efektywna}

\begin{lemat}[Okresowa przestrzeń Sobolewa]
\label{lem:pde30-periodic-H1}
Przestrzeń $H^1_{\mathrm{per}}(Y)$ jest domkniętą
podprzestrzenią $H^1(Y)$, a gładkie funkcje okresowe
są w niej gęste. Dla każdego $v\in H^1_{\mathrm{per}}(Y)$
zachodzi $\int_Y\nabla v\,\mathrm d\lambda_d=0$.
Podprzestrzeń $H^1_{\mathrm{per},0}(Y)$ jest przestrzenią
Hilberta z normą $\|\nabla v\|_{L^2(Y)}$.
\end{lemat}
\begin{proof}
Niech $v_n\in H^1_{\mathrm{per}}(Y)$ oraz
$v_n\to v$ w $H^1(Y)$.
Oznaczmy okresowe przedłużenia przez $\widetilde v_n$.
Dla dodatniej liczby całkowitej $m$ sześcian
$B_m=(-m,m)^d$ składa się, z dokładnością
do zbioru miary zero, z $(2m)^d$ komórek jednostkowych.
Okresowość funkcji i ich słabych pochodnych daje
\[
 \|\widetilde v_n-\widetilde v_l\|_{H^1(B_m)}^2
 =(2m)^d\|v_n-v_l\|_{H^1(Y)}^2.
\]
Ciąg jest Cauchy'ego w każdej z przestrzeni $H^1(B_m)$.
Z ich zupełności otrzymujemy granice zgodne na
przecięciach (przez ciągłość ograniczenia).
Definiują one funkcję $\widetilde v\in H^1_{\mathrm{loc}}(\R^d)$,
która ogranicza się do $v$ na $Y$.
Dla każdego $k\in\mathbb Z^d$ mamy
$\widetilde v(\cdot+k)=\widetilde v$ prawie wszędzie:
na dowolnym ograniczonym sześcianie obie strony
są granicami w $L^2$ odpowiednich przesunięć
$\widetilde v_n$. Zatem $v\in H^1_{\mathrm{per}}(Y)$.

Dla ustalonego $v\in H^1_{\mathrm{per}}(Y)$
wygładzamy $\widetilde v$ przez splot z gładkim
jądrem o zwartym nośniku dążącym do zera.
Sploty są okresowe, ponieważ przesunięcie o
$k\in\mathbb Z^d$ nie zmienia funkcji pod całką.
Zbieżność w $H^1(Y)$ wynika z lokalnego
twierdzenia o aproksymacji przez splot:
można pomnożyć $\widetilde v$ przez gładką
funkcję o zwartym nośniku równą $1$ w otoczeniu
$\overline Y$, gdyż dla małych jąder nie zmienia
to splotu na $Y$. Dla gładkich funkcji okresowych
całka każdej pochodnej po $Y$ znika, dzięki
twierdzeniu Newtona--Leibniza w odpowiedniej
współrzędnej i równości wartości na przeciwległych
ścianach. Przechodząc do granicy w $H^1(Y)$,
otrzymujemy $\int_Y\nabla v\,\mathrm d\lambda_d=0$.

Warunek średniej zero definiuje domkniętą
podprzestrzeń, gdyż $\lambda_d(Y)=1$ i całka
jest ciągła w $L^2(Y)$. Z nierówności
Poincarégo--Wirtingera na kostce
(Twierdzenie~\ref{thm:pde9-cube-Poincare})
mamy $\|v\|_2\le\sqrt d\|\nabla v\|_2$ dla
$v$ o średniej zero. Norma gradientu jest więc
równoważna normie $H^1$ i pochodzi od iloczynu
skalarnego $(v,w)\mapsto\int_Y\nabla v\cdot\nabla w\,
\mathrm d\lambda_d$. Domkniętość zapewnia jej zupełność.
\end{proof}

\begin{twierdzenie}[Istnienie i jednoznaczność poprawek komórkowych]\label{thm:pde30-cell}
Niech $A$ spełnia~\eqref{eq:pde30-ellipticity}.
Dla każdego $j\in\{1,\ldots,d\}$ istnieje dokładnie jedna
funkcja $\chi_j\in H^1_{\mathrm{per},0}(Y)$ taka, że
\begin{equation}\label{eq:pde30-cell}
\int_Y A(y)(e_j+\nabla_y\chi_j(y))\cdot\nabla_y\varphi(y)
\,\mathrm d\lambda_d(y)=0
\quad\text{dla każdego }\varphi\in H^1_{\mathrm{per},0}(Y).
\end{equation}
Definiujemy wektory strumienia $Q_j(y)=A(y)(e_j+\nabla\chi_j(y))$
i macierz stałą
\begin{equation}\label{eq:pde30-ahom}
 A^{\mathrm{hom}}e_j=\int_YQ_j(y)\,\mathrm d\lambda_d(y).
\end{equation}
Wtedy $\diver_yQ_j=0$ w sensie dystrybucji na torusie,
$A^{\mathrm{hom}}$ jest symetryczna i
$\lambda|\xi|^2\le\xi^{\mathsf T}A^{\mathrm{hom}}\xi
\le\Lambda|\xi|^2$ dla każdego $\xi\in\R^d$.
\end{twierdzenie}
\begin{proof}
Z Lematu~\ref{lem:pde30-periodic-H1}
przestrzeń $H^1_{\mathrm{per},0}(Y)$ z normą
gradientu jest przestrzenią Hilberta.
Dla $v,\varphi$ z tej przestrzeni forma
$b(v,\varphi)=\int_Y A\nabla v\cdot\nabla\varphi\,
\mathrm d\lambda_d$ spełnia
$|b(v,\varphi)|\le\Lambda\|\nabla v\|_2\|\nabla\varphi\|_2$
oraz $b(v,v)\ge\lambda\|\nabla v\|_2^2$.
Funkcjonał
$\varphi\mapsto-\int_YAe_j\cdot\nabla\varphi\,
\mathrm d\lambda_d$ ma normę co najwyżej $\Lambda$,
bo $\lambda_d(Y)=1$. Twierdzenie Laxa--Milgrama
daje jedyne $\chi_j$, a test $\chi_j$ daje
$\|\nabla\chi_j\|_2\le\Lambda/\lambda$.
W szczególności $Q_j\in L^2_{\mathrm{per}}(Y;\R^d)$.

Równanie~\eqref{eq:pde30-cell} zachodzi dla wszystkich
testów z $H^1_{\mathrm{per}}(Y)$, gdyż odjęcie średniej
nie zmienia gradientu. Stąd $\diver_yQ_j=0$
w sensie dystrybucji na torusie jednostkowym.
Zachodzi to również na $\R^d$ dla okresowego
przedłużenia: dla $\varphi\in C_c^\infty(\R^d)$
funkcja $\Phi(y)=\sum_{k\in\mathbb Z^d}\varphi(y+k)$
jest gładka i okresowa, a suma jest lokalnie skończona.
Z okresowości i rozbicia $\R^d$ na komórki wynika
\[
 \int_{\R^d}Q_j\cdot\nabla\varphi\,\mathrm d\lambda_d
 =\int_YQ_j\cdot\nabla\Phi\,\mathrm d\lambda_d=0.
\]

Dla $\xi\in\R^d$ niech $\chi_\xi=\sum_j\xi_j\chi_j$.
Z liniowości równań komórkowych
$\int_Y A(\xi+\nabla\chi_\xi)\cdot\nabla\varphi\,
\mathrm d\lambda_d=0$ dla każdego testu okresowego.
Z definicji macierzy efektywnej mamy
$A^{\mathrm{hom}}\xi=\int_Y A(\xi+\nabla\chi_\xi)\,
\mathrm d\lambda_d$.
Dodając do iloczynu skalarnego z $\xi$ zerową
całkę odpowiadającą testowi $\varphi=\chi_\xi$, otrzymujemy
\begin{equation}\label{eq:pde30-energycell}
 \xi^{\mathsf T}A^{\mathrm{hom}}\xi
 =\int_Y A(y)(\xi+\nabla\chi_\xi)\cdot
             (\xi+\nabla\chi_\xi)\,\mathrm d\lambda_d(y)
 =\min_{v\in H^1_{\mathrm{per},0}(Y)}
 \int_Y (\xi+\nabla v)^{\mathsf T}A(y)(\xi+\nabla v)\,\mathrm d\lambda_d(y).
\end{equation}
Aby uzasadnić drugą równość, oznaczmy energię
komórkową przez $\mathcal E_\xi(v)$.
Dla $v\in H^1_{\mathrm{per},0}(Y)$ rozwijamy kwadrat:
\[
 \mathcal E_\xi(v)-\mathcal E_\xi(\chi_\xi)
 =\int_Y A\nabla(v-\chi_\xi)\cdot\nabla(v-\chi_\xi)
 \,\mathrm d\lambda_d\ge0,
\]
gdyż wyraz mieszany znika na mocy
~\eqref{eq:pde30-cell} z testem $v-\chi_\xi$.
Równość wymusza $\nabla(v-\chi_\xi)=0$,
a warunek średniej zero i nierówność Poincarégo
dają $v=\chi_\xi$.
Ponieważ $\int_Y\nabla v\,\mathrm d\lambda_d(y)=0$,
mamy $\int_Y|\xi+\nabla v|^2\,\mathrm d\lambda_d(y)
=|\xi|^2+\int_Y|\nabla v|^2\,\mathrm d\lambda_d(y)\ge|\xi|^2$.
Stosując tę tożsamość dla $v=\chi_\xi$ i dolną
granicę eliptyczności $A$, otrzymujemy
$\xi^{\mathsf T}A^{\mathrm{hom}}\xi\ge\lambda|\xi|^2$.
Dla górnej granicy wybieramy $v=0$ w problemie
minimalizacyjnym, co daje
$\xi^{\mathsf T}A^{\mathrm{hom}}\xi\le
\int_Y\xi^{\mathsf T}A(y)\xi\,\mathrm d\lambda_d
\le\Lambda|\xi|^2$. Symetrię otrzymujemy z tożsamości
$e_i\cdot A^{\mathrm{hom}}e_j
=\int_Y(e_i+\nabla\chi_i)\cdot A(e_j+\nabla\chi_j)
\,\mathrm d\lambda_d(y)$, wynikającej z testu $\chi_i$
w zadaniu komórkowym. Prawa strona jest symetryczna
względem $i,j$, gdyż $A=A^{\mathsf T}$.
\end{proof}

\begin{wniosek}[Ograniczenia harmoniczne i arytmetyczne]
\label{cor:pde30-harmonic-arithmetic}
W założeniach Twierdzenia~\ref{thm:pde30-cell}
dla każdej $\xi\in\R^d$ zachodzi
\[
 \xi^{\mathsf T}\left(\int_Y A(y)^{-1}
 \,\mathrm d\lambda_d(y)\right)^{-1}\xi
 \le\xi^{\mathsf T}A^{\mathrm{hom}}\xi
 \le\xi^{\mathsf T}\left(\int_Y A(y)
 \,\mathrm d\lambda_d(y)\right)\xi .
\]
\end{wniosek}
\begin{proof}
Górna nierówność wynika z wyboru $v=0$
w~\eqref{eq:pde30-energycell}.
Niech $B=\int_Y A^{-1}\,\mathrm d\lambda_d$.
Macierz $A^{-1}$ jest mierzalna (odwracanie
jest ciągłe na zbiorze macierzy dodatnio określonych)
i spełnia oszacowania $\Lambda^{-1}I\le A^{-1}\le\lambda^{-1}I$
w sensie form kwadratowych. Zatem $B$ jest
symetryczna i dodatnio określona.
Dla dowolnego $v\in H^1_{\mathrm{per},0}(Y)$,
$z=\xi+\nabla v$ oraz stałego $\eta\in\R^d$
nieujemność formy daje
\[
 0\le\int_Y (z-A^{-1}\eta)^{\mathsf T}
 A(z-A^{-1}\eta)\,\mathrm d\lambda_d
 =\mathcal E_\xi(v)-2\eta\cdot\xi+\eta^{\mathsf T}B\eta.
\]
Użyliśmy $\int_Y z\,\mathrm d\lambda_d=\xi$.
Wybieramy $\eta=B^{-1}\xi$, a następnie
$v=\chi_\xi$. Z~\eqref{eq:pde30-energycell}
otrzymujemy dolną nierówność.
Dla $d=1$ dolna granica jest osiągana dzięki
poprawce ze Stwierdzenia~\ref{prop:pde30-corrector1d}.
\end{proof}

\begin{uwaga}[Co jest mikroskopowe, a co makroskopowe?]
Funkcja $\chi_j$ jest okresową poprawką komórkową i zależy
od zmiennej szybkiej $y$. Macierz $A^{\mathrm{hom}}$ jest stała
w zmiennej wolnej $x$. Współczynniki $A(x/\eps)$ nie muszą
zbiegać punktowo ani silnie do $A^{\mathrm{hom}}$; zbiega
rozwiązanie i odpowiednio rozumiany strumień. Nawet słaba
granica $A(x/\eps)$ jest na ogół inna niż $A^{\mathrm{hom}}$.
\end{uwaga}

\section{Twierdzenie o homogenizacji: pełny dowód przez oscylujące testy}
\begin{twierdzenie}[Homogenizacja okresowa przy mierzalnych współczynnikach]
\label{thm:pde30-homogenization}
Niech $d\ge1$, $\Omega\subset\R^d$ będzie ograniczonym
obszarem Lipschitzowskim, a $A$ spełnia~\eqref{eq:pde30-ellipticity}.
Dla $f\in H^{-1}(\Omega)$ niech $u_\eps\in H_0^1(\Omega)$
będzie jedynym słabym rozwiązaniem~\eqref{eq:pde30-micro}.
Niech $u_0\in H_0^1(\Omega)$ będzie jedynym słabym
rozwiązaniem zagadnienia
\begin{equation}\label{eq:pde30-macro}
 -\diver(A^{\mathrm{hom}}\nabla u_0)=f
 \quad\text{w }\Omega,\qquad u_0|_{\partial\Omega}=0.
\end{equation}
Wówczas dla $\eps\downarrow0$ zachodzą zbieżności
\[
 u_\eps\rightharpoonup u_0\text{ w }H_0^1(\Omega),\quad
 u_\eps\to u_0\text{ w }L^2(\Omega),\quad
 A(x/\eps)\nabla u_\eps\rightharpoonup
 A^{\mathrm{hom}}\nabla u_0\text{ w }L^2(\Omega;\R^d).
\]
W szczególności twierdzenie obejmuje obciążenia
$f\in L^2(\Omega)$, utożsamiane z funkcjonałami
$v\mapsto\int_\Omega fv\,\mathrm d\lambda_d$.
\end{twierdzenie}
\begin{proof}
\textbf{Krok 1: zwartość i strumień.}
Z oszacowania energetycznego
$\|\nabla u_\eps\|_2\le\lambda^{-1}\|f\|_{H^{-1}}$
oraz nierówności Poincarégo rodzina $(u_\eps)$
jest ograniczona w $H_0^1(\Omega)$.
Dla dowolnego ciągu dodatnich $\eps_k\to0$
z refleksywności przestrzeni Hilberta wybieramy
słabo zbieżny podciąg w $H_0^1$.
Z twierdzenia Rellicha~\ref{thm:pde10-Rellich-p}
po dalszym wydzieleniu $u_{\eps_k}\to u$ silnie
w $L^2$. Granice są zgodne, ponieważ zanurzenie
$H_0^1\hookrightarrow L^2$ jest ciągłe.
Strumienie $q_\eps=A(x/\eps)\nabla u_\eps$ spełniają
$\|q_\eps\|_2\le(\Lambda/\lambda)\|f\|_{H^{-1}}$;
po dalszym wydzieleniu $q_\eps\rightharpoonup q$ w $L^2$.
Dla każdego $\varphi\in H_0^1(\Omega)$
mamy $\int_\Omega q_\eps\cdot\nabla\varphi\,
\mathrm d\lambda_d=\langle f,\varphi\rangle$.
Przechodząc do granicy, otrzymujemy tę samą
tożsamość dla $q$, czyli $\diver q=-f$.
W dalszym ciągu piszemy $\eps$ zamiast indeksów podciągu.

\textbf{Krok 2: oscylujące funkcje testowe.}
Dla $i\in\{1,\ldots,d\}$ zdefiniujmy
$w_i^\eps(x)=x_i+\eps\chi_i(x/\eps)$.
Reguła łańcuchowa dla przeskalowania w $H^1_{\mathrm{loc}}$
daje $\nabla w_i^\eps=e_i+(\nabla\chi_i)(x/\eps)$.
Można ją uzasadnić przez gładkie okresowe przybliżenia
z Lematu~\ref{lem:pde30-periodic-H1}: przy ustalonym
$\eps$ podstawienie zmiennych i pokrycie komórkami
zapewniają zbieżność funkcji oraz gradientów w $L^2(\Omega)$.
W szczególności $w_i^\eps\in H^1(\Omega)$.
Mamy
$Q_i^\eps(x)=Q_i(x/\eps)=A(x/\eps)\nabla w_i^\eps(x)$.
Z Lematu~\ref{lem:pde30-average} oraz zerowej średniej
$\nabla\chi_i$ otrzymujemy
\[
 w_i^\eps\to x_i\ \text{w }L^2(\Omega),
 \qquad \nabla w_i^\eps\rightharpoonup e_i\
 \text{w }L^2(\Omega;\R^d).
\]
Zatem $w_i^\eps\rightharpoonup x_i$ w $H^1(\Omega)$.
Ten sam lemat daje
$Q_i^\eps\rightharpoonup A^{\mathrm{hom}}e_i$ w $L^2$.
Ponadto $\diver Q_i^\eps=0$ w sensie dystrybucji:
z dowodu Twierdzenia~\ref{thm:pde30-cell} wiemy,
że $\diver_yQ_i=0$ na całej $\R^d$.
Dla $\varphi\in C_c^\infty(\Omega)$ przedłużonej
przez zero podstawiamy $x=\eps y$:
\[
 \int_\Omega Q_i^\eps\cdot\nabla_x\varphi\,
 \mathrm d\lambda_d(x)
 =\eps^{d-1}\int_{\R^d} Q_i(y)\cdot
 \nabla_y(\varphi(\eps y))\,\mathrm d\lambda_d(y)=0.
\]
Nie wykorzystujemy klasycznych pochodnych $A$
ani ograniczoności $\chi_i$.

\textbf{Krok 3: identyfikacja granicy iloczynu.}
Dla $\varphi\in C_c^\infty(\Omega)$ iloczyn
$\varphi w_i^\eps$ należy do $H_0^1(\Omega)$:
ma zwarty nośnik, należy do $H^1$ na mocy reguły
iloczynu i można go wygładzić wewnątrz $\Omega$.
Ponadto mnożenie przez $\varphi$ jest ograniczonym
operatorem $H^1(\Omega)\to H_0^1(\Omega)$,
więc $\varphi w_i^\eps\rightharpoonup\varphi x_i$
w $H_0^1$. Testując równanie przez $\varphi w_i^\eps$
i rozwijając gradient iloczynu, otrzymujemy
\begin{align*}
\int_\Omega\varphi\,q_\eps\cdot\nabla w_i^\eps
\,\mathrm d\lambda_d(x)
&=\langle f,\varphi w_i^\eps\rangle
 -\int_\Omega w_i^\eps q_\eps\cdot\nabla\varphi
 \,\mathrm d\lambda_d(x)\\
&\longrightarrow \langle f,\varphi x_i\rangle
 -\int_\Omega x_i q\cdot\nabla\varphi\,\mathrm d\lambda_d(x)
 =\int_\Omega\varphi q_i\,\mathrm d\lambda_d(x).
\end{align*}
Pierwszy składnik zbiega dzięki słabej zbieżności
$\varphi w_i^\eps$ w $H_0^1$; w drugim odejmujemy
$\int_\Omega x_iq_\eps\cdot\nabla\varphi\,\mathrm d\lambda_d$,
a błąd jest ograniczony przez
$\|\nabla\varphi\|_\infty\|w_i^\eps-x_i\|_2\|q_\eps\|_2$.
Pozostała całka zbiega ze słabej zbieżności $q_\eps$.
Ostatnia równość wynika z testu $\varphi x_i$
w równaniu dla $q$.
Z drugiej strony, zerową dywergencję $Q_i^\eps$
możemy testować przez $\varphi u_\eps\in H_0^1$
dzięki gęstości funkcji testowych i $Q_i^\eps\in L^2$.
Otrzymujemy
\begin{align*}
\int_\Omega\varphi\,Q_i^\eps\cdot\nabla u_\eps
\,\mathrm d\lambda_d(x)
&=-\int_\Omega u_\eps Q_i^\eps\cdot\nabla\varphi
\,\mathrm d\lambda_d(x)\\
&\longrightarrow
 -\int_\Omega u(A^{\mathrm{hom}}e_i)\cdot\nabla\varphi
 \,\mathrm d\lambda_d(x)
 =\int_\Omega\varphi(A^{\mathrm{hom}}e_i)\cdot\nabla u
 \,\mathrm d\lambda_d(x).
\end{align*}
Tutaj błąd zastąpienia $u_\eps$ przez $u$ jest
ograniczony przez
$\|\nabla\varphi\|_\infty\|u_\eps-u\|_2\|Q_i^\eps\|_2$,
a pozostała całka zbiega ze słabej zbieżności
$Q_i^\eps$. Ostatnia równość to całkowanie przez
części dla stałego wektora $A^{\mathrm{hom}}e_i$.
Kluczowa jest dokładna tożsamość
$q_\eps\cdot\nabla w_i^\eps
=Q_i^\eps\cdot\nabla u_\eps$, wynikająca z symetrii $A$.
Porównując granice, dostajemy
$q_i=(A^{\mathrm{hom}}e_i)\cdot\nabla u$ dla każdego $i$.
Ponieważ $A^{\mathrm{hom}}$ jest symetryczna, oznacza to
$q=A^{\mathrm{hom}}\nabla u$.

\textbf{Krok 4: identyfikacja rozwiązania i cały ciąg.}
Równość $\diver q=-f$ i poprzednia identyfikacja dają
$-\diver(A^{\mathrm{hom}}\nabla u)=f$, $u\in H_0^1(\Omega)$.
Istnienie i jednoznaczność rozwiązania granicznego
wynikają z Twierdzenia~\ref{thm:pde30-cell}
i twierdzenia Laxa--Milgrama, zatem $u=u_0$.
Dla każdego ciągu $\eps_k\to0$ znaleźliśmy
podciąg realizujący wszystkie trzy zbieżności
do tych samych granic. Gdyby nie zachodziła
silna zbieżność w $L^2$, istniałby ciąg z normą
błędu ograniczoną od dołu przez dodatnią liczbę.
Analogicznie dla każdej z dwóch słabych zbieżności
testujemy przez ustalony ciągły funkcjonał liniowy.
W obu przypadkach otrzymalibyśmy sprzeczność
z własnością podciągu. Dowodzi to zbieżności całej
rodziny. Argument od początku obejmuje
$f\in H^{-1}(\Omega)$, bez aproksymowania $f$
funkcjami z $L^2$.
\end{proof}


\begin{twierdzenie}[Zbieżność energii i stabilność względem danych]
\label{thm:pde30-energy-convergence}
Niech $\Omega\subset\R^d$ będzie ograniczonym obszarem
Lipschitzowskim, a $A$ będzie symetryczna, okresowa,
mierzalna i jednostajnie eliptyczna jak w
\eqref{eq:pde30-ellipticity}. Dla $f\in H^{-1}(\Omega)$
niech $u_\eps$ i $u_0$ będą rozwiązaniami zagadnień
\eqref{eq:pde30-micro} i~\eqref{eq:pde30-macro}.
Wówczas
\begin{equation}\label{eq:pde30-energy-convergence}
 \lim_{\eps\downarrow0}\int_\Omega
 A(x/\eps)\nabla u_\eps\cdot\nabla u_\eps
 \,\mathrm d\lambda_d(x)
 =\int_\Omega A^{\mathrm{hom}}\nabla u_0\cdot\nabla u_0
 \,\mathrm d\lambda_d(x).
\end{equation}
Co więcej, jeśli $f_\eps\to f$ silnie w $H^{-1}(\Omega)$,
a $v_\eps\in H_0^1(\Omega)$ rozwiązuje
$-\diver(A(x/\eps)\nabla v_\eps)=f_\eps$,
to $v_\eps\rightharpoonup u_0$ w $H_0^1$,
$v_\eps\to u_0$ w $L^2$, a odpowiednie strumienie
zbiegają słabo w $L^2$. Zachodzi także
\eqref{eq:pde30-energy-convergence} po zastąpieniu
$u_\eps$ przez $v_\eps$.
\end{twierdzenie}
\begin{proof}
Testując równanie dla $u_\eps$ funkcją $u_\eps$,
otrzymujemy, przy oznaczeniu parowania dualnego,
\[
 \int_\Omega A(x/\eps)\nabla u_\eps\cdot\nabla u_\eps
 \,\mathrm d\lambda_d(x)=\langle f,u_\eps\rangle.
\]
Słaba zbieżność w $H_0^1$ z
\ref{thm:pde30-homogenization} oznacza zbieżność
prawej strony do $\langle f,u_0\rangle$.
Test równania granicznego funkcją $u_0$ identyfikuje
tę liczbę z prawą stroną~\eqref{eq:pde30-energy-convergence}.

Dla zmiennego obciążenia niech $u_\eps$ nadal oznacza
rozwiązanie przy ustalonym $f$. Różnica
$z_\eps=v_\eps-u_\eps$ spełnia równanie wariacyjne
$a_\eps(z_\eps,\varphi)=\langle f_\eps-f,\varphi\rangle$.
Test $\varphi=z_\eps$ i koercywność dają
$\lambda\|z_\eps\|_{H_0^1}^2\le
\|f_\eps-f\|_{H^{-1}}\|z_\eps\|_{H_0^1}$, czyli
\[
 \lambda\|z_\eps\|_{H_0^1}\le
 \|f_\eps-f\|_{H^{-1}},
\]
gdzie w normie $H_0^1$ bierzemy normę gradientu.
W szczególności $z_\eps\to0$ silnie w $H_0^1$,
a $\|A(x/\eps)\nabla z_\eps\|_{L^2}\to0$.
Wszystkie wymienione zbieżności dziedziczymy więc
z twierdzenia~\ref{thm:pde30-homogenization}.
Tożsamość energii dla $v_\eps$ ma postać
$\int_\Omega A(x/\eps)\nabla v_\eps\cdot\nabla v_\eps\,\mathrm d\lambda_d(x)
=\langle f_\eps,v_\eps\rangle$.
Z silnej zbieżności $f_\eps$ w $H^{-1}$,
ograniczoności $v_\eps$ w $H_0^1$ i jego słabej
zbieżności do $u_0$ wynika
$\langle f_\eps,v_\eps\rangle\to\langle f,u_0\rangle$:
różnica jest sumą
$\langle f_\eps-f,v_\eps\rangle+\langle f,v_\eps-u_0\rangle$.
Pierwszy składnik dąży do zera przez ograniczoność
norm $v_\eps$, drugi przez słabą zbieżność.
To kończy dowód.
\end{proof}

\begin{wniosek}[Lokalna zbieżność gęstości energii]
\label{cor:pde30-local-energy}
W założeniach Twierdzenia~\ref{thm:pde30-energy-convergence}
dla każdego $\varphi\in C_c(\Omega)$ zachodzi
\[
 \begin{aligned}
 &\lim_{\eps\downarrow0}\int_\Omega
 \varphi A(x/\eps)\nabla v_\eps\cdot\nabla v_\eps
 \,\mathrm d\lambda_d\\
 &\qquad=\int_\Omega
 \varphi A^{\mathrm{hom}}\nabla u_0\cdot\nabla u_0
 \,\mathrm d\lambda_d .
 \end{aligned}
\]
Obejmuje to przypadek $v_\eps=u_\eps$, $f_\eps=f$.
\end{wniosek}
\begin{proof}
Najpierw niech $\varphi\in C_c^\infty(\Omega)$.
Testując przez $\varphi v_\eps\in H_0^1$, mamy
\[
 \int_\Omega\varphi A(x/\eps)\nabla v_\eps\cdot\nabla v_\eps
 \,\mathrm d\lambda_d
 =\langle f_\eps,\varphi v_\eps\rangle
 -\int_\Omega v_\eps A(x/\eps)\nabla v_\eps\cdot\nabla\varphi
 \,\mathrm d\lambda_d .
\]
Mnożenie przez $\varphi$ jest ograniczone w $H_0^1$,
zatem pierwszy składnik zbiega do
$\langle f,\varphi u_0\rangle$ dzięki silnej zbieżności
$f_\eps$ i słabej zbieżności $\varphi v_\eps$.
Drugi składnik zbiega do
$-\int_\Omega u_0 A^{\mathrm{hom}}\nabla u_0\cdot\nabla\varphi
\,\mathrm d\lambda_d$ dzięki silnej zbieżności $v_\eps$
w $L^2$ i słabej zbieżności strumieni.
Stosując równanie graniczne z testem $\varphi u_0$,
identyfikujemy granicę z żądaną całką.
Dla $\varphi\in C_c(\Omega)$ wybieramy gładkie
przybliżenia zbieżne jednostajnie, z nośnikami
w ustalonym zbiorze zwartym wewnątrz $\Omega$.
Błąd aproksymacji jest ograniczony przez normę
jednostajną różnicy testów pomnożoną przez
$\Lambda\|\nabla v_\eps\|_2^2$, jednostajnie
ograniczoną. Analogicznie dla granicy.
Otrzymujemy tezę.
\end{proof}

\begin{uwaga}[Zbieżność energii nie oznacza silnej zbieżności gradientów]
W~\eqref{eq:pde30-energy-convergence} występują
\emph{zmienne formy kwadratowe}; nie jest to
zbieżność $\|\nabla u_\eps-\nabla u_0\|_2\to0$.
Przykład jednowymiarowy z
\ref{prop:pde30-corrector1d} pokazuje, że ta ostatnia
zbieżność może być fałszywa.
\end{uwaga}


\begin{uwaga}[Dlaczego w homogenizacji występuje energia?]
\label{rem:pde-v44-homogenization-monotone}
Dla każdego $\eps>0$ macierz $A(x/\eps)$
określa liniowy operator monotoniczny
$V\to V'$ z rozdziału~\ref{chap:pde26}. Istnienie
rozwiązania wynika z koercywności; nie wymaga ono
własnej konstrukcji energii homogenizacyjnej.
Jednak w twierdzeniu~\ref{thm:pde30-energy-convergence}
\emph{symetria} $A$ pozwala interpretować
$\frac12\int_\Omega A(x/\eps)\nabla v\cdot\nabla v\,
\mathrm d\lambda_d(x)-\langle f,v\rangle$ jako funkcjonał,
którego punkt minimalny rozwiązuje równanie. Dla
macierzy niesymetrycznych forma pozostaje koercywna,
lecz jej część antysymetryczna nie jest widoczna
w funkcjonale kwadratowym. Oszacowanie energii z testu
$v=u_\eps$ nadal ma sens; nie daje samo przez się
charakteryzacji minimizacyjnej całego operatora.
\end{uwaga}


\section{Warstwy w wymiarze dwa: anizotropia z izotropowych faz}
\begin{przyklad}[Efektywna macierz materiału warstwowego]\label{ex:pde30-layer}
Niech $d=2$, $A(y)=a(y_1)I_2$, gdzie $a$ jest okresowa,
$0<\lambda\le a\le\Lambda$. Dla $j=1$ poprawka zależy
tylko od $y_1$ i spełnia
$1+\chi_1'(y_1)=a_{\mathrm{hom}}/a(y_1)$.
Dla $j=2$ funkcja $\chi_2=0$ rozwiązuje zadanie komórkowe,
ponieważ $\diver_y(Ae_2)=\partial_{y_2}a(y_1)=0$.
Stąd
\begin{equation}\label{eq:pde30-layers}
 A^{\mathrm{hom}}=
 \begin{pmatrix}
  \bigl(\int_0^1 a(y)^{-1}\,\mathrm d\lambda_1(y)\bigr)^{-1}&0\\
  0&\int_0^1a(y)\,\mathrm d\lambda_1(y)
 \end{pmatrix}.
\end{equation}
Jeśli $a$ przyjmuje wartości $1$ i $4$ na dwóch równych
połówkach komórki, to $A^{\mathrm{hom}}=\operatorname{diag}(8/5,5/2)$.
W każdej fazie materiał jest lokalnie izotropowy,
ale jego zachowanie efektywne jest anizotropowe.
\end{przyklad}

\begin{stwierdzenie}[Granice średniej harmonicznej i arytmetycznej]
\label{prop:pde30-bounds}
Dla $d=1$ i współczynnika $a$ jak wyżej zachodzi
$\lambda\le a_{\mathrm{hom}}\le\int_0^1 a(y)
\,\mathrm d\lambda_1(y)\le\Lambda$,
a równość między średnią harmoniczną i arytmetyczną
zachodzi wtedy i tylko wtedy, gdy $a$ jest stała prawie wszędzie.
\end{stwierdzenie}
\begin{proof}
Pierwsza i ostatnia nierówność wynikają z oszacowań $a$.
Nierówność Cauchy'ego--Schwarza zastosowana do
$\sqrt a$ i $1/\sqrt a$ daje
$1\le(\int_0^1a)(\int_0^1a^{-1})$.
Równość zachodzi dokładnie, gdy $\sqrt a$ jest
proporcjonalna do $1/\sqrt a$ prawie wszędzie.
\end{proof}

\section{Pułapki zbieżności i poprawki pierwszego rzędu}
\begin{kontrprzyklad}[Słaba zbieżność współczynników nie określa równania granicznego]
Na $(0,1)$ niech $a(y)$ będzie niejednorodna i okresowa.
Wtedy $a(x/\eps)\rightharpoonup\int_0^1 a$ w $L^2(0,1)$,
lecz rozwiązania równań $-(a(x/\eps)u_\eps')'=f$
zbiegną do rozwiązania równania z $a_{\mathrm{hom}}
=(\int_0^1a^{-1})^{-1}$. Obie liczby różnią się,
gdy $a$ nie jest stała. Niedozwolone jest więc przejście
do granicy osobno w iloczynie dwóch słabo zbieżnych funkcji
$a(x/\eps)$ i $u_\eps'$.
\end{kontrprzyklad}

\begin{uwaga}[Granica zakresu twierdzenia]
Twierdzenie~\ref{thm:pde30-homogenization} dowodzi zbieżności
bez oszacowania jej szybkości w wymiarach $d\ge2$. Nie
implikuje samo z siebie globalnego szacowania
$\|\nabla u_\eps-(I+\nabla\chi(x/\eps))\nabla u_0\|_2\le C\eps$:
takie oszacowanie wymaga dodatkowej regularności danych,
współczynników i brzegu oraz analizy poprawek brzegowych.
Nie wystarcza formalne podstawienie ansatzu
$u_\eps\approx u_0+\eps\sum_j\chi_j(x/\eps)\partial_j u_0$.
W wymiarze jeden otrzymaliśmy natomiast
oszacowanie gradientowe z poprawką w
stwierdzeniu~\ref{prop:pde30-corrector1d} bez dodatkowej
gładkości współczynnika $a$.
\end{uwaga}

\section{Przykład numeryczny i interpretacja strumienia}
Dla $\Omega=(0,1)$ wybieramy $a(y)=2+\sin(2\pi y)$ i $f=1$.
Wtedy $\bar a=2$, podczas gdy
$a_{\mathrm{hom}}=\bigl(\int_0^1(2+\sin(2\pi y))^{-1}
\,\mathrm d\lambda_1(y)\bigr)^{-1}=\sqrt3$.
Rozwiązanie graniczne jest równe
$u_0(x)=x(1-x)/(2\sqrt3)$, a dokładne $u_\eps$ daje
wzór~\eqref{eq:pde30-explicit-1d}. W każdym mikrośrodowisku
strumień $q_\eps=a(x/\eps)u_\eps'$ jest funkcją afiniczną
$x$, mimo że sam gradient $u_\eps'$ szybko oscyluje.

W dodatku numerycznym~\ref{sec:pde30-code} przedstawiamy
kod Plotly, który wyznacza rodzinę rozwiązań $u_\eps$,
porównuje średnie arytmetyczną i harmoniczną, ilustruje
oscylacje gradientów oraz bada zbieżność numeryczną.



